
%------------------------------------------------------------------------------
%       package includes
%------------------------------------------------------------------------------
    % font encoding is set up for pdflatex, for other environments see
    % http://tex.stackexchange.com/questions/44694/fontenc-vs-inputenc
    \usepackage[T1]{fontenc}  % 8-bit fonts, improves handling of hyphenations
    % provides `old' commands for table of contents. Eases the ability to switch
    % between book and scrbook
    \usepackage{scrhack}


    % ------------------- layout, default -------------------
    % adjust the style of float's captions, separated from text to improve readabilty
    \usepackage[labelfont=bf, labelsep=colon, format=hang, textfont=singlespacing]{caption}
    % With format = hang your caption will look like this:
    % Figure 1: Lorem ipsum dolor sit amet,
    %           consectetuer adipiscing elit.
    %           Ut purus elit, vestibulum
    % If you instead want
    % Figure 1: Lorem ipsum dolor sit amet,
    % consectetuer adipiscing elit. Ut purus
    % elit, vestibulum
    % change to format=plain
    \usepackage{chngcntr}  % continuous numbering of figures/tables over chapters
    \counterwithout{equation}{chapter}
    \counterwithout{figure}{chapter}
    \counterwithout{table}{chapter}

    % Uncomment the following line if you switch from scrbook to book
    % and comment the setkomafont line
    \usepackage{titlesec}  % remove "Chapter" from the chapter title
    \titleformat{\chapter}[hang]{\bfseries\huge}{\thechapter}{2pc}{\huge}
    \titlespacing*{\chapter}{0pt}{-60pt}{20pt}
    % \setkomafont{chapter}{\normalfont\bfseries\huge}

    \usepackage{setspace}  % Line spacing
    \onehalfspacing
    % \doublespacing  % uncomment for double spacing, e.g. for annotations in correction

    % ------------------- functional, default-------------------
    % \usepackage{natbib}
    \usepackage[
        backend=biber,
        style=numeric,
        sorting=none,
        language=auto,
        autolang=other
    ]{biblatex}

    \addbibresource{bib/topic1.bib}
    \usepackage[dvipsnames]{xcolor}  % more colors
    \usepackage{array}  % custom format per column in table - needed on the title page
    \usepackage{graphicx}  % include graphics
    \usepackage{subfig}  % divide figure, e.g. 1(a), 1(b)...
    \usepackage{amsmath}  % |
    \usepackage{amsthm}   % | math, bmatrix etc
    \usepackage{amsfonts} % |
    \usepackage{calc}  % calculate within LaTeX
    \usepackage[unicode=true,bookmarks=true,bookmarksnumbered=true,
                bookmarksopen=true,bookmarksopenlevel=1,breaklinks=false,
                pdfborder={0 0 0},backref=false,colorlinks=false]{hyperref}


    %==========================================
    % You might not need the following packages, I only included them as they
    % are needed for the example floats
    % ------------------- functional, custom -------------------
    \usepackage{algorithm,algpseudocode}
    \usepackage{bm}  % bold greek variables (boldmath)
    \usepackage{tikz}
    \usetikzlibrary{positioning, arrows.meta}
    % \usetikzlibrary{positioning}  % use: above left of, etc

    % Improves general appearance of the text
    \usepackage[protrusion=true,expansion=true, kerning]{microtype}
    \usepackage{enumitem}

    % usually you don't need this, just for demonstration of a longer caption
    \usepackage{lipsum}

    \usepackage{parskip}

    % Own stuff:
    \usepackage[
        top=35mm, % original value: 43
        bottom=30mm,
        left=39mm, % original value: 43
        right=39mm % original value: 43
    ]{geometry}
    \usepackage[automark,headsepline]{scrlayer-scrpage}
    \clearpairofpagestyles
    \renewcommand{\chaptermark}[1]{%
        \markright{Kapitel \thechapter: #1}%
    }
    \ihead{\rightmark}
    \cfoot[\pagemark]{\pagemark}
    \renewcommand*{\chaptermarkformat}{}

    \usepackage[table]{xcolor}
    \usepackage{pifont}
    \usepackage{makecell}
    \usepackage{array}
    \usepackage{listings}
    \usepackage{pgfplots}
    \pgfplotsset{compat=1.18}
    % \usepackage{float}
    \newcolumntype{C}[1]{>{\centering\arraybackslash}p{#1}}

    \newcommand{\tabYes}{\cellcolor{green}\ding{51}}
    \newcommand{\tabPartial}{\cellcolor{blue!60}$\sim$}
    \newcommand{\tabNo}{\cellcolor{yellow}\ding{55}}
    \newcommand{\tabUnused}{$\sim$}
    \definecolor{codegray}{rgb}{0.9,0.9,0.9}
    \lstdefinestyle{mystyle}{
        backgroundcolor=\color{codegray},
        basicstyle=\ttfamily\small,
        breaklines=true,
        breakatwhitespace=true,
        columns=fullflexible,
        keepspaces=true,
        showstringspaces=false,
        showtabs=false,
        tabsize=2,
        frame=single,
        rulecolor=\color{gray},
        captionpos=b,
        xleftmargin=0.5em,
        xrightmargin=0.5em,
        framexleftmargin=0.5em,
        framexrightmargin=0.5em,
        aboveskip=1em,
        belowskip=1em,
        escapeinside={(*@}{@*)}
    }
    \lstset{style=mystyle}
    \renewcommand{\lstlistingname}{Codeblock}

    \newcommand{\benchmarkplot}[2]{%
    \begin{tikzpicture}
    \begin{axis}[
        width=0.8\textwidth,
        height=7cm,
        xlabel={Abdeckungsgrad der Testcases in archive in \%},
        ylabel={Zeit in ms},
        grid=major,
        scaled y ticks=false,
        yticklabel style={
            /pgf/number format/fixed,
            /pgf/number format/precision=0
        },
        legend style={
            at={(0.5,1.07)},
            anchor=south,
            legend columns=2
            }
        ]
        \addplot[
            thick,
            mark=*,
        ] coordinates {#1};
        \addlegendentry{konventioneller Testlauf}

        \addplot[
            thick,
            mark=square*,
            gray!100,
        ] coordinates {#2};
        \addlegendentry{Testlauf mit Verfahren}
    \end{axis}
    \end{tikzpicture}%
    }

    \newcommand{\benchmarkplotOverhead}[4]{%
    \begin{tikzpicture}
    \begin{axis}[
        width=0.8\textwidth,
        height=7cm,
        xlabel={Abdeckungsgrad der Testcases in archive in \%},
        ylabel={Zeit in ms},
        grid=major,
        scaled y ticks=false,
        yticklabel style={
            /pgf/number format/fixed,
            /pgf/number format/precision=0
        },
        legend style={
            at={(0.5,1.07)},
            anchor=south,
            legend columns=2
            }
        ]
        \addplot[
            thick,
            mark=*,
        ] coordinates {#1};
        \addlegendentry{testDiscovery}

        \addplot[
            thick,
            mark=square*,
            gray!140,
        ] coordinates {#2};
        \addlegendentry{testAuditor}

        \addplot[
            thick,
            mark=triangle*,
            gray!120,
        ] coordinates {#3};
        \addlegendentry{testExecution}

        \addplot[
            thick,
            mark=diamond*,
            gray!90,
        ] coordinates {#4};
        \addlegendentry{sonstiger Overhead}
    \end{axis}
    \end{tikzpicture}%
    }
%------------------------------------------------------------------------------
%       (re)new commands / settings
%------------------------------------------------------------------------------
    % ----------------- referencing ----------------
    \newcommand{\secref}[1]{Section~\ref{#1}}
    \renewcommand{\eqref}[1]{Equation~(\ref{#1})}
    \newcommand{\figref}[1]{\hyperref[#1]{Abbildung~\ref*{#1}}}
    \newcommand{\tabref}[1]{\hyperref[#1]{Tabelle~\ref*{#1}}}
    \newcommand{\chapref}[1]{\hyperref[#1]{Kapitel~\ref*{#1}}}
    \newcommand{\lstref}[1]{\hyperref[#1]{Codeblock~\ref*{#1}}}

    % ------------------- colors -------------------
    \definecolor{darkgreen}{rgb}{0.0, 0.5, 0.0}
    % Colors of the Albert Ludwigs University as in
    % https://www.zuv.uni-freiburg.de/service/cd/cd-manual/farbwelt
    \definecolor{UniBlue}{RGB}{0, 74, 153}
    \definecolor{UniRed}{RGB}{193, 0, 42}
    \definecolor{UniGrey}{RGB}{154, 155, 156}


    % ------------------- layout -------------------
    % prevents floating objects from being placed ahead of their section
    \let\mySection\section\renewcommand{\section}{\suppressfloats[t]\mySection}
    \let\mySubSection\subsection\renewcommand{\subsection}{\suppressfloats[t]\mySubSection}


    % ------------------- marker commands -------------------
    % ToDo command
    \newcommand{\todo}[1]{\textbf{\textcolor{red}{(TODO: #1)}}}
    \newcommand{\extend}[1]{\textbf{\textcolor{darkgreen}{(EXTEND: #1)}}}
    % Lighter color to note down quick drafts
    \newcommand{\draft}[1]{\textbf{\textcolor{NavyBlue}{(DRAFT: #1)}}}


    % ------------------- math formatting commands -------------------
    % define vectors to be bold instead of using an arrow
    \renewcommand{\vec}[1]{\mathbf{#1}}
    \newcommand{\mat}[1]{\mathbf{#1}}
    % tag equation with name
    \newcommand{\eqname}[1]{\tag*{#1}}


    % ------------------- pdf settings -------------------
    % ADAPT THIS
    \hypersetup{pdftitle={\thetitle},
                pdfauthor={\theauthor},
                pdfsubject={Undergraduate thesis at the Albert Ludwig University of Freiburg},
                pdfkeywords={junit xml, testing, nix},
                pdfpagelayout=OneColumn, pdfnewwindow=true, pdfstartview=XYZ, plainpages=false}


    %==========================================
    % You might not need the following commands, I only included them as they
    % are needed for the example floats

    % ------------------- Tikz styles -------------------
    \tikzset{>=latex}  % arrow style


    % ------------------- algorithm ---------------------
    % Command to align comments in algorithm
    \newcommand{\alignedComment}[1]{\Comment{\parbox[t]{.35\linewidth}{#1}}}
    % define a foreach command in algorithms
    \algnewcommand\algorithmicforeach{\textbf{foreach}}
    \algdef{S}[FOR]{ForEach}[1]{\algorithmicforeach\ #1\ \algorithmicdo}

    % line spacing should be 1.5
    \renewcommand{\baselinestretch}{1.5}

    % set distance between items in a list, for more details see the
    % enumitem package: https://www.ctan.org/pkg/enumitem
    \setlist{itemsep=.5em}
